\chapter{Conclusion and Future Scope}

\section{Summary of Project Work}
This project addressed the critical clinical and computational challenges of transitioning modern radiation oncology toward an \textbf{MRI-Only Radiotherapy Planning Workflow}. By synthesizing accurate, geometrically faithful synthetic Computed Tomography (sCT) volumes directly from magnetic resonance imaging (MRI), the need for planning CT simulation scans can be eliminated—eradicating inter-modality coregistration misalignments ($1\text{--}5\text{ mm}$), preventing unnecessary ionizing radiation burdens in cancer patients, and significantly streamlining clinical workflows.

Formulated on the international SynthRAD2023 Grand Challenge cranial cohort (comprising 60 clinical brain patients and 3,614 paired 2D axial slices), this work implemented, trained, and benchmarked five distinct generative deep learning architectures under strictly identical patient-level splits ($\text{Seed} = 42$, $80\% / 20\%$) and standardized evaluation protocols:
\begin{enumerate}
    \item \textbf{Model 1: Deep Residual U-Net (ResUNet Baseline)}—Supervised convolutional regression baseline optimized via composite masked $L_1$ and Mean Squared Error losses.
    \item \textbf{Model 2: Conditional GAN (Pix2Pix)}—Paired adversarial translation framework featuring a U-Net Generator and a $70 \times 70$ PatchGAN Discriminator.
    \item \textbf{Model 3: Conditional DDPM From Scratch}—Likelihood-based generative diffusion model trained directly on paired MRI-CT distributions.
    \item \textbf{Model 4: SOTA-CDDPM with Residual Attention}—Conditional diffusion model enhanced with cross-layer self-attention gating and an anatomical bone-weighted hybrid objective ($\mathbf{w}_{\text{bone}} = 5.0$).
    \item \textbf{Model 5: Proposed Cross-Conditioned Residual Diffusion (CCR-Diffusion)}—A novel two-stage decoupled architecture featuring a deterministic Stage 1 ResUNet coarse anatomical anchor coupled with a Stage 2 diffusion denoiser operating on standardized residuals dynamically modulated via Feature-wise Linear Modulation (FiLM).
\end{enumerate}

\section{Key Findings and Benchmark Insights}
Comprehensive quantitative benchmarking and hardware profiling across the 363 held-out test slices established several decisive scientific and clinical findings:
\begin{itemize}
    \item \textbf{CCR-Diffusion Achieves Landmark Performance:} The proposed CCR-Diffusion architecture emerged as the definitive winner across all four quantitative clinical metrics, attaining an overall $\text{MAE}_{\text{all}}$ of $\mathbf{48.2 \pm 14.6\text{ HU}}$, a bone-specific $\text{MAE}_{\text{bone}}$ of $\mathbf{98.4 \pm 19.2\text{ HU}}$, a Structural Similarity Index ($\text{SSIM}$) of $\mathbf{0.9248 \pm 0.0215}$, a Peak Signal-to-Noise Ratio ($\text{PSNR}$) of $\mathbf{26.85 \pm 1.94\text{ dB}}$, and a Bone Dice coefficient of $\mathbf{0.8245 \pm 0.0482}$.
    \item \textbf{Resolution of the Regression-Blur Dilemma:} Standard supervised regressors (Model 1) converge to the conditional expectation $\mathbb{E}[\mathbf{y} \mid \mathbf{x}]$, leading to oversmoothed bone boundaries and elevated bone errors ($168.4\text{ HU}$). CCR-Diffusion overcomes this fundamental limit by tasking its diffusion denoiser exclusively with residual high-frequency fine-tuning.
    \item \textbf{Elimination of Adversarial Hallucination Risks:} While conditional GANs (Model 2) produce visually sharp images, they are prone to mode collapse and fictitious pseudo-calcifications that pose grave risks of beam underdosing in radiotherapy. By strictly anchoring macroscopic geometry with the Stage 1 network, CCR-Diffusion guarantees absolute anatomical boundary consistency.
    \item \textbf{Resolution of Air vs. Bone Hypointensity Ambiguity:} In raw $T_1$-weighted MRI, cortical bone ($T_2^* < 0.5\text{ ms}$) and air cavities both present as dark, signal-void regions. While baseline models frequently misclassified paranasal sinuses as intermediate bone, CCR-Diffusion successfully disambiguated air ($-1000\text{ HU}$) from petrous bone ($> 1200\text{ HU}$) by leveraging multi-scale structural priors from the anchor encoder.
    \item \textbf{Clinical Real-Time Feasibility:} Utilizing a 50-step deterministic Denoising Diffusion Implicit Model (DDIM, $\eta = 0$) reverse sampling schedule, CCR-Diffusion achieves an inference throughput of approximately $2.01\text{ FPS}$ ($498\text{ ms}$ per slice). A full 3D cranial radiotherapy volume of 120 slices can be synthesized in under 60 seconds on a standard clinical GPU accelerator, fitting comfortably within routine radiation oncology simulation workflows. In contrast to iterative multi-atlas non-rigid registration pipelines that routinely require 30 to 120 minutes per patient, this sub-minute generative turnaround empowers same-session online adaptive replanning directly at the linear accelerator console. Furthermore, the deterministic nature of zero-variance DDIM trajectory caching guarantees bitwise-reproducible attenuation maps across repeated clinical executions, satisfying rigorous radiation therapy quality assurance (QA) and medical device verification protocols.
\end{itemize}

\section{Limitations of the Current Work}
Despite its significant advancements, the current study possesses several limitations that warrant consideration:
\begin{enumerate}
    \item \textbf{Two-Dimensional In-Plane Slice Processing:} The current framework operates on independent 2D axial slices. Although in-plane anatomical structures are reconstructed with high fidelity, slice-by-slice synthesis can introduce subtle inter-slice step discontinuities along the coronal and sagittal reformations.
    \item \textbf{Computational Latency vs. Pure Feedforward Models:} Although 50-step DDIM reduces sampling time by $20\times$ compared to 1,000-step DDPM, an inference latency of $498\text{ ms}$ per slice is still substantially slower than single-pass feedforward networks ($18.4\text{ ms}$). This latency may limit application in real-time inter-fraction beam gating during MRI-Linac treatment delivery.
    \item \textbf{Provisional HU Calibration Window:} Experiments were conducted using normalized dynamic ranges mapped to the SynthRAD2023 cranial window $[-1000\text{ HU}, +1500\text{ HU}]$. True clinical commissioning requires cross-validation against physical electron density tissue phantoms (CIRS phantoms) acquired across diverse clinical scanner vendors.
\end{enumerate}

\section{Future Research Scope}
The promising results achieved by CCR-Diffusion open several compelling avenues for future technical innovation and clinical translation:
\begin{enumerate}
    \item \textbf{Three-Dimensional Volumetric and Latent Diffusion Models (LDMs):} Extending CCR-Diffusion to 3D convolutional or Latent Diffusion topologies to operate natively on 3D spatial volumes, guaranteeing seamless volumetric continuity across all orthogonal planes (axial, sagittal, coronal).
    \item \textbf{Few-Step Consistency Models and Rectified Flow:} Investigating Consistency Distillation (CD) and Rectified Flow algorithms to compress the diffusion reverse trajectory from 50 steps down to 1--4 steps, enabling ultra-fast real-time inference ($< 20\text{ ms}$ per slice) for online adaptive MRI-guided radiotherapy on MRI-Linac systems.
    \item \textbf{End-to-End Dosimetric Validation and Gamma Analysis:} Ingesting the generated sCT volumes directly into commercial treatment planning systems (e.g., Eclipse, RayStation, Monaco) to compute full 3D Monte Carlo radiation dose distributions. Evaluating clinical gamma passing rates ($\gamma_{3\%/3\text{mm}} > 98\%$ and $\gamma_{2\%/2\text{mm}} > 95\%$) and Dose-Volume Histograms (DVHs) on Planning Target Volumes (PTVs) and critical Organs-at-Risk (brainstem, optic nerves, and chiasm).
    \item \textbf{Multi-Contrast MRI Ingestion:} Extending the conditioning framework to ingest complementary MRI sequences ($T_1$-weighted, $T_2$-weighted, and FLAIR) to provide rich, orthogonal tissue contrast for enhanced anatomical delineation.
    \item \textbf{Cross-Anatomical and Multi-Center Generalization:} Validating the CCR-Diffusion framework on other challenging anatomical sites—including pelvic prostate cohorts (SynthRAD2023 Task 2) and head-and-neck cancer patients—to establish a unified, generalizable clinical synthetic CT engine.
\end{enumerate}

\section{Concluding Remarks}
In conclusion, this project successfully formulated, implemented, and benchmarked five deep learning paradigms for MRI-to-CT translation, demonstrating the decisive superiority of the proposed \textbf{CCR-Diffusion} framework. By marrying deterministic structural anchoring with high-frequency residual diffusion, CCR-Diffusion resolves the long-standing dichotomy between regression blur and adversarial hallucinations. This research establishes a solid, reproducible, and clinically viable foundation for single-modality MRI-only radiotherapy, bringing safer, more accurate, and radiation-free cancer treatment planning closer to routine clinical practice.
